
This section compares the \ptbar\ spectra between the 
  \PbPb\ and \pp\ data, to gain some understanding of how the
  spectra changes from the \pp\ to peripheral \PbPb\ to central \PbPb\ collisions.
Figures~\ref{fig:ptbar_ch2}-~\ref{fig:ptbar_ch1}
  show ratios of the \ptbar\ spectra for All, Same-sign
  and Opposite-sign pairs, respectively.
The plotted spectra are for the difference between the Away-Side 
  and Near-side pairs, i.e. the difference in the \ptbar\ spectra 
  for pairs with $|\Dphi|>\pi/2$ and pairs with $|\Dphi|<\pi/2$
  (and not for the yields obtained from the fits Eq.~\ref{eq:fit_func}--\ref{eq:signal_func}).
The plots demonstrate that the most central \PbPb\ collisions (0--10\%)
  have a relatively larger contributions from muons at higher \ptbar\
  and a deplection at lower \ptbar\ compared to the \pp\ collisions. 
This might partially contribute to the effect seen in Figure~\ref{fig:result_widths},
  where the widths of the \Dphi\ distributons in most-central collisions are narrower.
%This is presemably because the quenching of particles/hadrons in the QGP
%  decreases with increasing \pt\ leading to a stronger depletion at lower pt.

Figure~\ref{fig:ptbar_tight_4GeV} shows the mean-\ptbar\ values, 
  as a function of centrality 
  (for the difference of distributions on the away-side and near-side)
  for the case of \ptbar>4~\GeV.
Figure~\ref{fig:ptbar_tight} shows the corresponding plot
  for the nominal $\ptbar>5~\GeV$ case.
In both cases, the mean values are calculated with the restriction 
  that $\ptbar<20~\GeV$.
For the \ptbar>4~\GeV\ case, it is observed that for peripheral events (60--80\% centrality),
  the \pp\ and \PbPb\ $\langle\ptbar\rangle$ values are in agreement.
For more central collisions the $\langle\ptbar\rangle$ values
  for the \PbPb\ data are higher, especially for the opposite-sign case,
  indicating a stronger modification of the \pt\ spectra, 
  in most central collisions.
While for the \ptbar>5~\GeV\ case, the \PbPb\ values for all centralitiies
  are compareable with the \pp\ values.
This is one of the reasons that the \ptbar>5~\GeV\ requirement
  was used in the measurement, as over this \ptbar\ region
  the spectra are more comparable between the \pp\ and \PbPb\ data.

\begin{figure}[h]
\centering
\includegraphics[width=1.0\textwidth]%
{figures/AllFigs/can_pTBar_option1_ch2_dphi3_Tight_pT_GT4_EffCor}
%{figures/AllFigs/can_pTBar_option1_ch2_dphi3_Tight_pTBar_GT5_EffCor}
\caption{
Ratio of the single-muon spectra in \PbPb\ collisions to that
  in \pp\ collisions.
The plotted spectra are for the difference between the Away-Side 
  and Near-side pairs,
The plots have an overall arbitrary scale (different for each panel).
Each panel corresponds to a different centrality interval.
The muons are from All-sign pairs.
\label{fig:ptbar_ch2}
}
\end{figure}

\begin{figure}[h]
\centering
\includegraphics[width=1.0\textwidth]%
{figures/AllFigs/can_pTBar_option1_ch0_dphi3_Tight_pT_GT4_EffCor}
%{figures/AllFigs/can_pTBar_option1_ch0_dphi3_Tight_pTBar_GT5_EffCor}
\caption{
Same as Figure~\ref{fig:ptbar_ch2} but for muons from Same-sign pairs.
\label{fig:ptbar_ch0}
}
\end{figure}

\begin{figure}[h]
\centering
\includegraphics[width=1.0\textwidth]%
{figures/AllFigs/can_pTBar_option1_ch1_dphi3_Tight_pT_GT4_EffCor}
%{figures/AllFigs/can_pTBar_option1_ch1_dphi3_Tight_pTBar_GT5_EffCor}
\caption{
Same as Figure~\ref{fig:ptbar_ch2} but for muons from Opposite-sign pairs.
\label{fig:ptbar_ch1}
}
\end{figure}


\begin{figure}[h] 
\centering
\includegraphics[width=1.0\textwidth]%
{figures/AllFigs/can_pTBar_option1_Tight_pT_GT4_EffCor}
%\includegraphics[width=1.0\textwidth]%
%{figures/AllFigs/can_pTBar_option1_pT_GT4_EffCor}
\caption{
The$\langle\ptbar\rangle$ values for the muon-pairs
  contributing to the pair distributions.
The mean values are calculated for the difference of the $\ptbar$
  distributions on the away-side ($|\Dphi|<\pi2$) 
  and near-side ($|\Dphi|>\pi/2$), with the restriction
  that $4<\ptbar/\GeV<20$.
The \PbPb\ values are shown for different centrality intervals ($x$-axis).
The horizontal line indicates the value for the \pp\ data,
  which are shown for multiple cuts on the minimum \ptbar\
  to get an idea of how different the \PbPb\ and \pp\ spectra are.
The left, center and right panels correspond to same-charge,
  opposite-charge and combined-charge pairs.
%The top and bottom panels correspond to Tight and Medium muons, respectively.
\label{fig:ptbar_tight_4GeV}
}
\end{figure}

\begin{figure}[h] 
\centering
\includegraphics[width=1.0\textwidth]%
{figures/AllFigs/can_pTBar_option1_Tight_pTBar_GT5_EffCor}
%\includegraphics[width=1.0\textwidth]%
%{figures/AllFigs/can_pTBar_option1_pTBar_GT5_EffCor}
\caption{
Same as Figure~\ref{fig:ptbar_tight_4GeV} but with the requirement
  that \ptbar>5~\GeV.
\label{fig:ptbar_tight}
}
\end{figure}



\clearpage
\section{\ptbar\ spectra plots for paper}
\label{sec:meanpTv2}
Figures~\ref{fig:ptbar_ch2_v2}--\ref{fig:ptbar_ch1_v2}
  show ratios of the \ptbar\ spectra for All, Same-sign
  and Opposite-sign pairs, respectively.
The plotted spectra are for the difference between the Away-Side 
  and Near-side pairs,
These are similar to those shown in Figures~\ref{fig:ptbar_ch2}--\ref{fig:ptbar_ch1_v2}
  but are plotted over the range $\ptbar>5~\GeV$ and includes systematic uncertainties
  that arise from 
  1) Trigger and Reco Efficiencies,
  2) Muon working point,
  3) The $|\Delta\eta|$ cut
     and
  4) Drell-Yan corrections.
%Overall normalization uncertainties are dropped out due to the plots having 
%  an arbritary scale.
 
\begin{figure}[h]
\centering
\includegraphics[width=1.0\textwidth]%
{figures/AllFigs/can_pTBar_dphi3_ch2_Tight_pTBar_GT5_EffCor}
\caption{
Ratio of the single-muon spectra in \PbPb\ collisions to that
  in \pp\ collisions.
The plotted spectra are for the difference between the Away-Side 
  and Near-side pairs,
The plots have an overall arbitrary scale (different for each panel).
Each panel corresponds to a different centrality interval.
The muons are from All-sign pairs.
This figure is similar to Figure~\ref{fig:ptbar_ch2},
  but includes systematic uncertainties and shows the spectra only for $\ptbar>5~\GeV$.
\label{fig:ptbar_ch2_v2}
}
\end{figure}

\begin{figure}[h]
\centering
\includegraphics[width=1.0\textwidth]%
{figures/AllFigs/can_pTBar_dphi3_ch0_Tight_pTBar_GT5_EffCor}
\caption{
Same as Figure~\ref{fig:ptbar_ch2_v2} but for muons from Same-sign pairs.
\label{fig:ptbar_ch0_v2}
}
\end{figure}

\begin{figure}[h]
\centering
\includegraphics[width=1.0\textwidth]%
{figures/AllFigs/can_pTBar_dphi3_ch1_Tight_pTBar_GT5_EffCor}
\caption{
Same as Figure~\ref{fig:ptbar_ch2_v2} but for muons from Opposite-sign pairs.
\label{fig:ptbar_ch1_v2}
}
\end{figure}
